% Copyright 2026 Alan Szmyt. All rights reserved pending publication license.
\section{Methods}
\label{sec:methods}

This paper specifies a prospective design and feasibility protocol rather than
a clinical efficacy study. The canonical machine-readable artifact is
\texttt{ANT-PROT-FEAS-001}, version~1.0.0, frozen on 2026-09-07 with status
\texttt{frozen-design-protocol}. Its merge grants no collection authority: no
formal human study has begun, all human-collection activation gates remain
unsatisfied, and no ethics or safety approval is presumed. The design draws on
N-of-1 protocol guidance for prospective specification and complete reporting,
pilot and feasibility guidance for separating practicability from efficacy,
and a digital N-of-1 platform as an implementation precedent
\cite{porcino2020spent,eldridge2016feasibility,konigorski2022studyu}. Those
sources inform reporting discipline; they do not validate Antidote's measures,
conditions, software, or proposed use.

\subsection{Planned measurement layers}
\label{sec:planned-measurement-layers}

The protocol preserves intended semantic controls, model and adapter identity,
generation settings, output identity and checksum, measured acoustic features,
actual exposure, baseline and target states, separated subjective responses,
timestamps, protocol version, and mutations between trials. These layers are
non-interchangeable. A requested control is not a realized acoustic property; a
measured acoustic property is not perceived expression; perceived expression
is not felt state; and emotional intensity, completion, or engagement is not
evidence of usefulness, harm, or clinical benefit. Behavioral or physiological
measures are optional extensions and cannot become the sole definition of a
person's experienced state.

\subsection{Evidence classification}
\label{sec:evidence-classification}

Research records use five non-interchangeable classes: \textbf{source} for what
external evidence establishes, \textbf{hypothesis} for what the project proposes,
\textbf{observation} for what occurred, \textbf{interpretation} for a bounded
explanation of an observation, and \textbf{claim} for a statement supported in
the manuscript. The working claim ledger is maintained outside publication
artifacts under \texttt{research/notes/}. Synthetic fixture output and repository
tests remain technical observations; they are not participant data.

\subsection{Research questions and feasibility stages}
\label{sec:methods-feasibility-stages}

The protocol maps the three research questions to four evidence stages. D0
addresses RQ1 through design traceability: each transformation and authority
boundary must resolve to a versioned record, equation, contract, or explicit
unavailable state. T0 addresses the implemented portion of RQ2 with synthetic,
deterministic conformance and recovery tests. T1 would extend RQ2 to a frozen
real-model, adapter, analyzer, hardware, and mutation manifest, but is currently
\texttt{blocked-no-real-model}. H1 would address RQ3 through a prospective
repeated-session within-person series, but is
\texttt{blocked-no-collection-authority}. Passing D0 or T0 neither implies that
T1 will pass nor authorizes H1. Likewise, technical controllability cannot
establish that an exposure is subjectively fitting, useful, safe, or
therapeutically effective.

\subsection{Units, conditions, and session structure}
\label{sec:methods-session-structure}

The technical unit is one immutable generation specification and its terminal
result; the human-response unit is one assigned session attempt, with an
exposure counted only when a verified artifact is actually played. The proposed
H1 series contains six blocks and three conditions per block, for 18 planned
completed exposures and an absolute ceiling of 24 scheduled attempts. Each
artifact targets eight minutes. Exposures are separated by at least 48 hours,
limited to three in any seven days, and followed by an aftereffect window
targeting 24 hours (acceptable range 18--30 hours).

Condition G is a generic instrumental ambient generation that receives neither
current state, desired transition, nor personal context. Condition S receives
only selected current-state and desired-transition categories and a frozen
meet--hold--transition--integrate template; it receives no free-text history,
personal descriptors, or prior-session mappings. Condition P receives the
explicitly approved working projection, person-authored semantic mixins,
inclusions, exclusions, and an editable journey plan. Across conditions, a
qualifying execution would hold constant the immutable model revision, adapter
and analyzer versions, hardware profile, duration, output format, playback
setup where feasible, response-instrument version, and file target of
$-23$~LUFS integrated loudness with a $-1$~dBTP true-peak ceiling. File-level
normalization neither determines sound-pressure level at the ear nor establishes
hearing safety; playback begins at the lowest device level and remains under
the person's control. Auditory beats, within-exposure adaptation, and personal-
model updating are excluded from this protocol.

Each prospective session follows eight governed steps: authority and readiness;
baseline and desired transition; condition reveal and plan review; generation
and verification; a separate, explicit playback action; immediate response
within five minutes; later aftereffect report; and closure or pause. A generated
but unplayed artifact never enters an exposure or human-response denominator.

\AntidoteFigure{feasibility-protocol-timeline}

\subsection{Technical verification}
\label{sec:methods-technical-verification}

T0 verifies the mock system's authority, schema, replay, interruption,
recovery, hashing, and provenance invariants with complete denominators. T1,
if activated after implementation review, freezes a real-model qualification
manifest before any qualifying run. Its planned matrix contains three baseline
runs, 24 one-factor mutation runs, six adjacent-segment continuity runs, and
six injected failure-and-recovery runs: 39 attempts if all declared controls
are supported. The fixed seeds are 101, 307, and 911.

The one-factor pairs vary tempo (70 versus 100~BPM), normalized density (0.25
versus 0.65), timbral instruction, or semantic intent while holding seed, model,
adapter, hardware, duration, output format, base plan, and all other controls
constant. Tempo and density require versioned measured outputs. Timbral
comparison requires a named analyzer and uncertainty. Semantic mutation is
verified only as an adapter-conditioning trace unless an independently
specified measure exists; it is not relabeled as acoustic adherence. Unsupported
controls remain unsupported rather than being silently substituted.

The technical measures are authority-state conformance, required provenance
completeness, component latency and verified buffer margin, requested-versus-
measured control adherence, semantic and acoustic boundary distances, and
terminal generation state. They correspond to the consent projection,
adapter, generate-ahead, continuity, and response-record equations
(ANT-EQ-003, ANT-EQ-007, and ANT-EQ-010--ANT-EQ-014). Hard gates require all
invalid or unapproved transitions to fail closed; every playable artifact to
match its declared hash, format, and bounded path; cancellation, timeout, stop,
and safety halt to reach recorded terminal states; and no missed deadline,
partial output, or unsupported control to be reported as success. Hardware- and
adapter-specific latency, buffer, analyzer, and perceptual thresholds must be
frozen in the later T1 manifest; this paper invents no universal threshold.

\subsection{Within-person feasibility measures}
\label{sec:methods-within-person-measures}

The proposed H1 instrument uses brief, near-event self-report because ecological
momentary assessment can reduce recall distance while introducing its own
burden, reactivity, adherence, and missingness concerns
\cite{shiffman2008ema}. Baseline captures a person-authored description plus
valence, arousal, and intensity; desired transition remains a separate category
and optional description. The compact valence and arousal interaction borrows
from the Affective Slider \cite{betella2016slider}, but Antidote's wording,
intensity, response timing, and additional items constitute a new protocol and
are not validated by that source.

After exposure, the instrument records perceived musical expression separately
from felt state. It then records whether the intensity was wanted, as well as
helpfulness, resonance, mismatch, harm, surprise, interaction burden, and total
session burden. Helpfulness and harm are asked again in the later aftereffect
window without overwriting the immediate record. Every bounded item uses its
declared anchors. A supplied felt-state object requires a description, while a
declined or unavailable felt-state response is null and named in the field-level
missingness record. The record also preserves the instrument version,
correction lineage, window, and missingness reason. No clinical
outcome is collected, and no emotional direction is prespecified as inherently
better. In particular, catharsis, crying, arousal, or movement toward a target
cannot substitute for the person's appraisal of whether the experience was
wanted, helpful, mismatched, or harmful.

Physiology is disabled in version~1.0.0. A future supplemental module would
require separate consent; device, calibration, clock, unit, signal-quality,
artifact, dropout, and missingness records; a frozen analysis; and independent
privacy and safety review. It could not overwrite self-report, diagnose emotion
or need, prove entrainment, or steer generation under this protocol.

\AntidoteTable{measurement-evidence-matrix}

\subsection{Assignment, carryover, and missingness}
\label{sec:methods-assignment-missingness}

The six possible orders of G, S, and P are each used once. Their block order is
randomized with a cryptographically generated seed committed in a separate
activation record before session~1. Eligibility and baseline are confirmed
before condition reveal. Participant blinding is impossible because the three
interfaces disclose different context and plan operations; researcher blinding
is also impossible in a self-study, and both facts must be reported. A failed
assigned session remains in the attempt ledger; a replacement may use the same
condition at the next eligible time but never erases the failure.

At least 48 hours and completion or expiry of the later-response window separate
exposures. Ongoing aftereffect, mismatch, distress, hearing discomfort, or
uncertainty postpones the next attempt by at least another 24 hours and triggers
a repeated readiness review. This spacing may reduce but cannot prove removal
of carryover. Flagged sessions remain in primary accounting and are examined in
a prespecified sensitivity analysis.

Analyses distinguish all assigned attempts, verified artifacts, started
exposures, completed exposures, immediate windows, and later windows. No outcome
is imputed in the primary analysis. Missing, declined, missed-window,
interrupted, and technical-failure states remain separate; condition contrasts
use only blocks containing both required observations and disclose every omitted
block. Best- and worst-case bounded values are used only as sensitivity checks
for feasibility progression thresholds, never as substituted observations.

\subsection{Analysis and sensitivity plan}
\label{sec:methods-analysis-plan}

Technical reporting includes complete state denominators, invariant outcomes,
median, range, and empirical latency quantiles, paired control differences by
factor and seed, and every semantic and acoustic boundary component. Human-
response feasibility reporting includes assigned through withdrawn counts by
condition, field-level completeness and reasons, burden distributions, all
mismatch, harm, hearing-discomfort, distress, and safety events, and the ability
to link semantic, acoustic, exposure, and response records through provenance.

Secondary within-person estimands are blockwise P-minus-S, P-minus-G, and
S-minus-G contrasts for immediate helpfulness, resonance, mismatch, harm, and
burden; differences between perceived-expression and felt-state valence and
arousal; descriptive rank associations among varied semantic controls, measured
features, and responses when at least six non-missing, nonconstant pairs exist;
and condition-specific baseline-to-immediate changes when both observations
exist. Every session point, denominator, median, range, and blockwise contrast
is shown. If all six blocks are analyzable, the analysis enumerates all
$6^6=46{,}656$ within-block condition-label permutations as an exploratory
within-person randomization distribution, not a population estimate or
confirmatory test.

Sensitivity analyses compare started exposures with those reaching 80 percent
of audio, include and exclude carryover flags, compare early and late blocks,
adjust descriptive models for block, order, and baseline state, leave out one
block at a time, identify self-researcher versus independently administered
sessions, distinguish fixed-seed from non-deterministic adapters, and compare
no-imputation results with bounded missingness scenarios. All prespecified null,
zero, opposite-direction, negative, missing, interrupted, burdensome, harmful,
and adverse observations remain visible. No endpoint will be selected or
renamed after observation.

The progression rule is operational rather than therapeutic. A later bounded
study may be considered only if the hard technical gates pass; at least 15 of
18 exposures and at least four per condition complete; at least 80 percent of
started exposures have an immediate record or explicit decline reason and at
least 70 percent have the same for later response; median interaction burden is
no greater than 0.5; no safety event remains unresolved; and the participant
and an independent reviewer separately find another bounded study acceptable.
Otherwise the decision is stop or revise. Even a ``proceed'' decision establishes
neither efficacy, generalizability, privacy, nor safety.

Equation ANT-EQ-014 defines the separated response record used here.
ANT-EQ-015 and ANT-EQ-016 remain unavailable future empirical models: this
protocol fits no response predictor, specifies no likelihood, and authorizes no
personal-model update.

\subsection{Protocol freeze, safety, and deviations}
\label{sec:methods-protocol-freeze}

The canonical protocol, its SHA-256 lock, and an initially empty append-only
deviation log are stored under \texttt{experiments/protocols/}. Before the first
qualifying record, a material change requires a new semantic protocol version,
content lock, analysis plan, and review. After collection begins under separate
authority, the frozen file is never rewritten: amendments and deviations append
the rationale, affected records, evidence impact, author, reviewer, and
disposition. Conditions, assignment, timing, measures, estimands, exclusions,
model or analyzer identity, consent, privacy, retention, safety, physiology,
and missingness rules are material changes.

The person may postpone, pause, skip, stop, withdraw, or decline any response
without penalty or automatic substitution. A stop request or consent revocation,
hearing pain or unexpected loudness, unwanted escalating distress or
disorientation, technical or provenance failure, or any loss of freely chosen
participation stops playback and prevents automatic continuation. The interface
must offer silence and a clear exit, retain only the minimum consented event,
and defer optional questions until the person chooses. Urgent concerns belong
with appropriate human or emergency support, not Antidote. Unresolved adverse
responses, serious events, stop-control failure, or repeated high harm pause the
protocol for independent review; these rules do not prove clinical safety.

Separate consent scopes govern inspection, projection, generation, playback,
retention, analysis, export, and any future personal-model use. Raw journal,
therapy-chat, health, or other private free text cannot enter the public
repository or a default research export. Storage encryption and recovery,
retention expiry and deletion, projection invalidation, backups, participant
materials, the applicable ethics determination, assignment commitment,
instrument implementation, analysis code, and pilot stop/go review all remain
unresolved activation gates. Consequently, this protocol is reproducible as a
design artifact but is not permission to recruit, expose, collect, or report a
human result.
